\section{Discussion}
This AJA-targeted framework makes three deliberate commitments. First, it selects KNHANES as the primary public dataset because it is Korean, population-based, and relevant to the intended clinical setting, with NHANES for cross-national external validation. Second, it reframes AI as a reproducible research accelerator rather than a diagnostic authority, bounded by a deterministic safety layer. Third, it operationalizes reproducibility and shared benefit through fully open code, a data dictionary, model cards, and a synthetic-data path that lets anyone run the pipeline without restricted data.

\subsection{Interpretation}
The expected scientific contribution is \emph{not} a claim that workplace stress directly damages hearing. Rather, StressEar-AI quantifies whether stress-related measures are consistently associated with tinnitus, hearing difficulty, and audiometric outcomes after accounting for noise exposure and health covariates. We anticipate a robust stress--tinnitus association and a weaker, noise-confounded stress--threshold association (H1). This asymmetry, if confirmed, has a clear clinical reading: stress may amplify tinnitus perception and burden more than it shifts cochlear thresholds, which aligns with tinnitus being partly a central, attention-modulated phenomenon.

\subsection{Clinical contribution and the treatment window}
The clinical contribution is an audiology-centered pathway that separates screening, referral, treatment, and rehabilitation. The framework's most consequential clinical stance is procedural rather than statistical: because SSNHL is a time-critical emergency \citep{chandrasekhar2019ssnhl}, the system is engineered so that a low predicted risk can never suppress urgent referral. The motivating vignette---resolved unilateral stress-associated tinnitus followed a year later by a delayed-care sudden loss in the other ear---illustrates precisely the failure mode the red-flag layer exists to prevent. Public data can identify population patterns and high-risk profiles; the Ajou University Hospital extension can then test treatment timing, SSNHL recovery, speech-in-noise outcomes, hearing-aid rehabilitation, and tinnitus burden longitudinally.

\subsection{Comparison with prior work}
Unlike single-country prevalence studies \citep{park2014tinnitus,joo2015hrqol}, StressEar-AI prespecifies cross-national transportability testing and reports calibration and decision-curve utility, not discrimination alone. Unlike opaque clinical-AI demonstrations, it locks feature and outcome definitions before external validation, publishes model cards, and evaluates subgroup fairness as a primary result. This positions the framework as a reusable methodological baseline rather than a one-off result.

\subsection{Limitations}
Several limitations remain and temper interpretation. Stress measurement differs across datasets and may not capture validated occupational-stress constructs; we therefore label the exposure ``perceived stress'' when appropriate. Cross-sectional audiometry cannot establish temporal order, so no causal claim is made. Audiometry availability and item wording vary across survey cycles, constraining pooling. Tinnitus definitions and severity scales differ across sources. Public datasets lack detailed treatment timing, which only the prospective cohort can supply. Complex-survey design must be honored to avoid biased inference. Finally, open medical language models require local validation, privacy review, and clinical governance before any deployment; their outputs are confined here to research feature extraction and guideline-linked explanation, never to autonomous clinical decisions.

\subsection{Future directions}
Future work includes: adding longitudinal occupational cohorts to strengthen temporality; a neurophysiological substudy of tinnitus subtypes using open OpenNeuro/EEG datasets; audiogram-digitization to unlock legacy paper audiograms for retrospective analysis; and community extension of the harmonization schema to additional national health surveys, advancing a globally shared, reproducible evidence base.
